

\section{Cơ sở xử lý tín hiệu cảm biến}

Trong các hệ thống giám sát tại biên, tiền xử lý và trích xuất đặc trưng là bước quan trọng ảnh hưởng trực tiếp đến hiệu suất của mô hình. Thay vì đẩy trực tiếp dữ liệu thô với số chiều lớn vào mạng nơ-ron, quy trình trích xuất đặc trưng giúp giữ lại các thông tin cốt lõi, giảm nhiễu và hạ thấp đáng kể yêu cầu tính toán. Đối với bài toán giám sát máy giặt, hệ thống khai thác đồng thời hai nguồn dữ liệu: tín hiệu âm thanh để nhận diện biến động tần số cao và tín hiệu rung động để theo dõi trạng thái cơ khí ở tần số thấp.

\subsection{Trích xuất đặc trưng âm thanh}
\label{subsec:audio_features}

Tín hiệu âm thanh thu từ microphone INMP441 qua bus I2S có dạng chuỗi thời gian liên tục. Để mô hình tự mã hóa nhận diện được các âm thanh bất thường như lỗi ma sát hay tiếng rít động cơ, kỹ thuật năng lượng dải lọc Mel được áp dụng. Phương pháp này mô phỏng thính giác con người — nhạy cảm với các thay đổi ở dải tần số thấp hơn là tần số cao. Phép biến đổi từ tần số Hz sang thang đo Mel được tính theo Phương trình \ref{eq:mel_scale}:

\begin{equation}
\label{eq:mel_scale}
M(f) = 2595 \cdot \log_{10} \left( 1 + \frac{f}{700} \right)
\end{equation}

\begin{figure}[H]
\centering
\resizebox{\textwidth}{!}{%
\begin{tikzpicture}[
    node distance=0.55cm,
    block/.style={draw, rounded corners=2pt, align=center, minimum height=0.9cm, text width=2.25cm, font=\small},
    arrow/.style={->, thick}
]
\node[block] (raw) at (0,0) {Tín hiệu I2S\\âm thanh};
\node[block] (frame) at (2.8,0) {Chia khung\\512 mẫu};
\node[block] (hann) at (5.6,0) {Cửa sổ\\Hann};
\node[block] (fft) at (8.4,0) {FFT\\$|X(k)|^2$};
\node[block] (mel) at (11.2,0) {13 bộ lọc\\Mel};
\node[block] (log) at (14.0,0) {Log-energy\\MFE};
\draw[arrow] (raw) -- (frame);
\draw[arrow] (frame) -- (hann);
\draw[arrow] (hann) -- (fft);
\draw[arrow] (fft) -- (mel);
\draw[arrow] (mel) -- (log);
\end{tikzpicture}
}
\caption{Đường ống trích xuất đặc trưng âm thanh MFE.}
\label{fig:audio_mfe_pipeline}
\end{figure}

Quy trình xử lý trong Hình \ref{fig:audio_mfe_pipeline} bắt đầu bằng việc chia tín hiệu thành các khung dài 64 ms, tương đương 512 mẫu ở tần số lấy mẫu 8 kHz, trượt với bước nhảy 32 ms. Các tham số này được lựa chọn dựa trên sự cân bằng giữa khả năng giữ lại thông tin âm học và giới hạn tài nguyên của vi điều khiển. Tần số lấy mẫu 8 kHz cho phép giữ lại dải âm thanh 0--4 kHz theo định lý Nyquist, đủ bao phủ các thành phần tiếng ồn cơ khí chính của động cơ, vòng bi và lồng giặt, trong khi giảm đáng kể dung lượng bộ đệm và số phép tính so với các mức 16 kHz hoặc 44,1 kHz \cite{ref_shul_2023, ref_1_16}. Với 8 kHz, khung 512 mẫu tương ứng 64 ms và đồng thời là kích thước lũy thừa của 2, phù hợp để thực hiện FFT hiệu quả. Cấu hình này tạo độ phân giải tần số khoảng $8000/512 \approx 15{,}6$ Hz, đủ chi tiết để quan sát sự thay đổi phổ của tiếng ma sát hoặc tiếng rít cơ khí; nếu khung ngắn hơn, độ phân giải tần số giảm, còn nếu khung dài hơn, độ trễ và bộ nhớ đệm tăng không phù hợp với xử lý tại biên.

Bước nhảy 32 ms tạo mức chồng lấn 50\% giữa hai khung liên tiếp. Mức chồng lấn này giúp đặc trưng phổ thay đổi mượt hơn theo thời gian, hạn chế bỏ sót các bất thường ngắn nhưng vẫn không làm tăng gấp đôi tải tính toán như các mức chồng lấn dày hơn. Trước khi biến đổi, cửa sổ Hann được nhân với từng khung để làm mềm hai biên của đoạn tín hiệu, qua đó giảm hiện tượng rò rỉ phổ khi FFT giả định tín hiệu trong mỗi khung có tính tuần hoàn \cite{ref_1_16}. Nhờ vậy, các đặc trưng năng lượng Mel ổn định hơn và ít bị ảnh hưởng bởi nhiễu biên khung.

Sau đó, phép biến đổi Fourier nhanh (FFT) chuyển tín hiệu sang miền tần số, và phần năng lượng phổ được đưa qua dải 13 bộ lọc hình tam giác xếp chồng theo thang Mel \cite{ref_1_12}. Số lượng 13 bộ lọc được lựa chọn nhằm cân bằng giữa khả năng phân giải phổ và kích thước vectơ đầu vào. Giá trị này đủ để mô tả các vùng năng lượng chính của tiếng ồn cơ khí, đồng thời giữ không gian đặc trưng âm thanh ở mức nhỏ để khi kết hợp với 6 đặc trưng rung động, vectơ đầu vào cuối cùng chỉ có 19 chiều và phù hợp với vùng nhớ hạn chế của TensorFlow Lite Micro trên vi điều khiển.

Đầu ra của mỗi bộ lọc là một giá trị năng lượng logarit. Với bộ lọc Mel thứ $m$, đặc trưng năng lượng được tính theo Phương trình \ref{eq:mel_energy}:

\begin{equation}
\label{eq:mel_energy}
E_m = \log \left( \sum_{k=0}^{K-1} |X(k)|^2 H_m(k) + \epsilon \right)
\end{equation}

Trong đó, $X(k)$ là phổ sau FFT, $H_m(k)$ là đáp ứng của bộ lọc Mel thứ $m$, $K$ là số điểm tần số được xét và $\epsilon$ là hằng số nhỏ để tránh phép tính $\log(0)$. Đồ án chủ đích chọn MFE thay vì MFCC (Mel-Frequency Cepstral Coefficients) nhằm loại bỏ bước biến đổi Cosin ngược (DCT). Lựa chọn này giúp tiết kiệm chu kỳ CPU trên vi điều khiển trong khi vẫn bảo toàn đủ thông tin không gian phổ để phát hiện bất thường \cite{ref_1_1, ref_1_12}.

\subsection{Trích xuất đặc trưng rung động cơ học}

Tín hiệu rung động từ cảm biến gia tốc ADXL345 phản ánh trực tiếp các sự cố cơ học như mất cân bằng tải hay lệch trục lồng giặt. Vì các lỗi này thường bộc lộ rõ qua cường độ và độ phân tán của dao động trong miền thời gian, đồ án tiến hành trích xuất hai đặc trưng thống kê là Giá trị hiệu dụng (RMS) và Phương sai (Variance) cho từng trục (X, Y, Z).

Giá trị RMS đo lường mức năng lượng trung bình của rung động, được tính theo Phương trình \ref{eq:rms}:

\begin{equation}
\label{eq:rms}
x_{RMS} = \sqrt{\frac{1}{N} \sum_{i=1}^{N} x_i^2}
\end{equation}

Bên cạnh đó, Phương sai đánh giá mức độ biến thiên của tín hiệu quanh giá trị trung bình. Việc kết hợp phương sai giúp mô hình nhạy hơn với các xung va đập bất thường do mất cân bằng tải gây ra \cite{ref_1_3}. Công thức tính phương sai được mô tả tại Phương trình \ref{eq:variance}:

\begin{equation}
\label{eq:variance}
\sigma^2 = \frac{1}{N} \sum_{i=1}^{N} (x_i - \bar{x})^2
\end{equation}

Chỉ sử dụng các đặc trưng miền thời gian giúp hệ thống không phải tính toán các phép biến đổi miền phức tạp (như FFT), qua đó rút ngắn tối đa độ trễ suy luận và tiết kiệm dung lượng SRAM trên vi điều khiển Seeed Studio XIAO ESP32-S3 \cite{ref_1_11}.

\subsection{Kỹ thuật kết hợp dữ liệu đa phương thức}

Để đánh giá đầy đủ trạng thái vận hành của thiết bị, đồ án áp dụng kỹ thuật kết hợp dữ liệu ở mức đặc trưng. Phương pháp này giúp mô hình bắt được mối tương quan chéo giữa âm thanh và rung động; chẳng hạn, tiếng ồn tần số cao đi kèm với biên độ rung lớn thường là dấu hiệu rõ ràng của lỗi ở pha vắt.

Vectơ đầu vào $V_{input} \in \mathbb{R}^{19}$ được tạo ra bằng cách ghép các đặc trưng theo Phương trình \ref{eq:fusion}:

\begin{equation}
\label{eq:fusion}
V_{input} = [MFE_{1 \dots 13} \oplus RMS_{x,y,z} \oplus \sigma^2_{x,y,z}]
\end{equation}

Trong đó, toán tử $\oplus$ biểu diễn phép ghép nối, gộp không gian 13 chiều MFE và 6 chiều rung động thành một vectơ đặc trưng 19 chiều duy nhất.

Thao tác quan trọng nhất ở bước này là sử dụng \textbf{Chuẩn hóa hỗn hợp} để tinh chỉnh từng loại dữ liệu trước khi nạp vào mạng tự mã hóa. Với đặc trưng âm thanh MFE, do cường độ luôn bị giới hạn vật lý bởi độ nhạy của microphone, hệ thống dùng \textbf{Min-Max Scaling} để ánh xạ dải $[-80, 0]$ dB về khoảng $[0, 1]$. Ngược lại, dữ liệu rung động (RMS và phương sai) thường xuyên xuất hiện các đỉnh nhiễu cao đột biến khi có va đập. Để xử lý, hệ thống dùng \textbf{Robust Scaling} — phương pháp chuẩn hóa dựa trên trung vị và dải phân vị. Cách này đưa dữ liệu về cùng thang đo chuẩn mà không làm lu mờ thông tin của các đỉnh giá trị bất thường \cite{ref_1_8}.

Đường ống tiền xử lý và chuẩn hóa hỗn hợp này không chỉ cung cấp dữ liệu sạch cho mô hình trí tuệ nhân tạo mà còn là nền tảng để hệ thống chạy đa nhân song song ổn định trên FreeRTOS. Các nghiên cứu liên quan và những hạn chế mà hệ thống này khắc phục sẽ được tổng kết ở cuối chương.

% --- Ghi chú Tài liệu tham khảo ---
% \bibitem{ref_1_1} P. P. Ray, "A review on TinyML...", 2022.
% \bibitem{ref_1_3} H. Mohammadi, "EARLY DETECTION OF IMBALANCE IN LOAD...", 2020.
% \bibitem{ref_1_8} M. Abdallah et al., "Anomaly Detection and Inter-Sensor Transfer Learning on Smart Manufacturing Datasets," Sensors, 2023.
% \bibitem{ref_1_11} T. S. Ajani et al., "An Overview of Machine Learning within Embedded and Mobile Devices--Optimizations and Applications," Sensors, 2021.
% \bibitem{ref_1_12} U. Muthumala et al., "Comparison of Tiny Machine Learning Techniques for Embedded Acoustic Emission Analysis," IEEE WF-IoT, 2024.
